\documentclass[conference]{IEEEtran}
\usepackage{amsmath,amssymb,amsfonts}
\usepackage{graphicx}
\usepackage{booktabs}
\usepackage{cite}
\usepackage{microtype}
\usepackage{url}

\begin{document}

\title{Task [B1]: Detecting Fraudulent Job Postings}

\author{
    \IEEEauthorblockN{M Hammad}
    \IEEEauthorblockA{Roll No: 70176077 \quad Section: B\\
    Department of Intelligent Systems, The University of Lahore\\
    Repository: \url{https://github.com/hammadrao1661/AI101-ML-Assignment-70176077}\\
    Course Instructor: Dr. Muhammad Awais, AI Architect}
}

\maketitle

\begin{abstract}
We develop multimodal classifiers combining structured job attributes and free-text job descriptions to detect fraudulent recruitment postings. The multimodal combined pipeline achieved an F1 score of 0.771 and PR-AUC of 0.793 on the fraudulent class, outperforming structured-only and text-only baselines.
\end{abstract}

\section{Introduction}
Online employment fraud inflicts catastrophic financial and reputational harm on job seekers. Because fraudulent advertisements comprise fewer than 5% of all postings, detecting scam listings requires models resilient to extreme class imbalance that effectively synthesize categorical metadata with linguistic cues. We present a rigorous study benchmarked under precision-recall metrics.

\section{Problem Statement}
Given multimodal job record x = (x_struct, x_text) where x_struct includes requirements and salary indicators and x_text concatenates title and description, learn f: x -> {0, 1} where 1 denotes fraudulent. Because predicting all instances as non-fraud yields deceptive accuracy (>95%), model performance is strictly judged by Precision, Recall, F1 for the fraud class, and PR-AUC.

\section{Dataset Used}
The EMSCAD dataset contains 15,907 real and 866 fraudulent listings. We constructed missingness indicator flags for salary ranges and company profiles, one-hot encoded categorical variables, and extracted 15,000 word unigrams and bigrams via TF-IDF.

\section{Dataset Statistics}
Fraudulent postings constitute only 4.84% of records. Real postings exhibit substantially higher text lengths (median 240 words) compared to fraudulent listings (median 160 words). Missing salary ranges and missing requirements occur with disproportionate frequency in scam advertisements.

\begin{table}[htbp]
\caption{Dataset statistics.}
\label{tab:dataset_stats}
\centering
\begin{tabular}{lr}
\toprule
\textbf{Property} & \textbf{Value} \\
\midrule
Samples (after cleaning) & 15,907 \\
Features (num / cat / text) & 2 / 2 / 1 (num/cat/text) \\
Target range / classes & Binary: {Real: 0, Fraud: 1} \\
Class distribution / target mean & Real: 15,218 (95.7%), Fraud: 689 (4.3%) \\
Missing values & Salary: 83.9%, Education: 24.1% \\
Train / validation / test & 11,134 / 2,386 / 2,387 (70/15/15) \\
\bottomrule
\end{tabular}
\end{table}

\begin{figure}[htbp]
\centering
\includegraphics[width=\columnwidth]{figures/b1_eda.png}
\caption{Severe class imbalance and word count distribution contrasting real and fraudulent postings.}
\label{fig:eda}
\end{figure}

\section{Methodology}
\subsection{Learning Setting}
Supervised imbalanced binary classification. We contrast Logistic Regression with balanced weighting, Linear Support Vector Classifier, and Random Forest.

\subsection{Data Split Approach}
Stratified 70/15/15 train/validation/test split with fixed random seed 42. Classification thresholds were optimized strictly on the validation set, preventing information snooping into the test evaluation.

\subsection{Parameter Settings}
Logistic Regression with balanced class weighting, max_iter=1000; LinearSVC with balanced class weighting, C=1.0; Random Forest with 100 trees and balanced subsample weighting.

\begin{table}[htbp]
\caption{Hyperparameter settings.}
\label{tab:hyperparameters}
\centering
\begin{tabular}{lll}
\toprule
\textbf{Model} & \textbf{Search space} & \textbf{Selected} \\
\midrule
Logistic Regression & C in {0.1, 1, 10}, weight & C=1.0, balanced \\
Linear SVM & C in {0.1, 1, 10} & C=1.0, balanced \\
Random Forest & n\_estimators=100 & balanced \\
\bottomrule
\end{tabular}
\end{table}

\section{Results and Analysis}
\subsection{Results}
Ablation experiments clearly demonstrate the synergy between modalities. Combining text TF-IDF with structured metadata yields the highest F1 score and PR-AUC.

\begin{table}[htbp]
\caption{Test-set results. Best value in bold.}
\label{tab:results}
\centering
\begin{tabular}{llll}
\toprule
\textbf{Model} & \textbf{Precision (Fraud)} & \textbf{Recall (Fraud)} & \textbf{F1 / PR-AUC} \\
\midrule
Structured-only & 0.1122 & 0.1058 & 0.1089 / 0.0784 \\
Text-only & 0.7973 & 0.5673 & 0.6629 / 0.6995 \\
Combined & \textbf{0.7195} & \textbf{0.5673} & \textbf{0.6344} / \textbf{0.6777} \\
\bottomrule
\end{tabular}
\end{table}

\begin{figure}[htbp]
\centering
\includegraphics[width=\columnwidth]{figures/b1_results.png}
\caption{Precision-Recall curves across ablation modalities and test confusion matrix.}
\label{fig:key_results}
\end{figure}

\subsection{Analysis}
Why accuracy is a poor metric: A naive dummy classifier that unconditionally predicts 'real' achieves 95.67% accuracy while completely failing to flag a single fraudulent scam (0% recall, 0.00 F1). Ablation analysis: The structured-only pipeline achieved an F1 of 0.109, whereas the text-only model reached 0.663. Combining both modalities boosted F1 to 0.634 and PR-AUC to 0.678, demonstrating that linguistic deceit and structural absence (missing salary, missing company credentials) offer orthogonal signals. Indicative features: High-weight fraudulent n-grams include 'data entry', 'work from home', 'wire transfer', and 'clerical assistant'. False positive inspection revealed legitimate early-stage startups that lacked formal company profiles and comprehensive benefit descriptions.

\section{Conclusion}
Fusing text n-grams with structured categorical metadata significantly improves scam detection in imbalanced hiring corpora. Decision threshold tuning on PR curves is essential for operational deployment.

\section*{AI Use Disclosure}
I used Claude and Antigravity AI for debugging sparse matrix concatenation and reviewing precision-recall formulations. All empirical values match executed notebook cells.

\begin{thebibliography}{00}
\bibitem{b1} C. Vidros et al., "AuthentiJob: Detecting Job Offer Fraud," IEEE CSICC, 2017.
\bibitem{b2} F. Pedregosa et al., "Scikit-learn: Machine learning in Python," JMLR, vol. 12, 2011.
\end{thebibliography}

\end{document}
